% ============================================================
\section{Results}
\label{sec:results}
% ============================================================

\subsection{Behavioural Results}
\label{sec:behavioural-results}

We report on the 23 models that are not already reproducibly correct on the
untreated question. The fleet is 24 distinct models measured under two protocols:
seventeen sampled through hosted APIs and nine open-weight models sampled locally in
bf16 with a per-sample seed recorded, two of which appear in both and are counted
once from the local rows. All models receive byte-identical prompt text at $T=1.0$,
$\mathrm{top}_p = 1.0$, with $R=10$ samples per condition. One model (Qwen2.5-3B)
answers correctly on all ten untreated samples and is excluded throughout: a model
that is already reproducibly correct cannot exhibit a rescue, and including it would
let an aggregate figure move without any model changing state. The remaining 23 give
230 executions per condition.

Results are reported as model counts by operational state. Aggregate correctness
appears only where it is the object of the argument rather than the measure of the
outcome.

\paragraph{Untreated operational state.}
Under the untreated question, sixteen of the 23 models are reproducibly incorrect,
naming the wrong action on all ten samples. The remaining seven are
non-reproducible, returning both actions across repetitions of one unchanged input.
None is reproducibly correct, by construction of the study population.

The two states fail differently and only one of them fails visibly. A
non-reproducible model announces its own unreliability: sample it twice and the
answers disagree. A reproducibly incorrect model does not. For sixteen of
twenty-three models --- roughly seven in ten --- the failure survives any procedure
that samples the same input repeatedly and compares, because every sample agrees.
Repetition is the standard check available to a practitioner at deployment time, and
on this decision it certifies the majority of the fleet as stable while they are
uniformly wrong.

\paragraph{Effects of conventional prompt interventions.}
Seven conventional interventions were applied to the same question under the same
protocol: expert role assignment, objective emphasis, chain-of-thought elicitation,
anti-hallucination instruction, threat, encouragement, and error-avoidance
instruction. None changes the operational state of the fleet
(Table~\ref{tab:conditions}). The strongest, expert role assignment, converts one of
the sixteen reproducibly incorrect models to reproducible correctness. No
conventional condition converts more than one, and across all seven the converted
model is the same one each time. Five of the seven leave at least fifteen of the
sixteen reproducibly incorrect.

Four of the seven also turn a non-reproducible model into a reproducibly incorrect
one: one each under chain-of-thought, anti-hallucination, threat and error
avoidance, and two under encouragement. Such a model is both less accurate and
harder to diagnose, since its ten samples now agree. Under encouragement the number
of reproducibly incorrect models is therefore higher than untreated, at seventeen of
twenty-three.

\paragraph{Specification-conditioned operational state.}
Under the task specification, supplied as a system message with the question
unchanged, seventeen of the 23 models are reproducibly correct. Two remain
reproducibly incorrect --- Qwen2.5-7B and Qwen3-4B-2507, whose internal measurements
in Section~\ref{sec:internal-results} show large but insufficient displacement
rather than no effect --- and four are non-reproducible.

No model becomes reproducibly incorrect under the specification, a transition that
four of the seven conventional interventions produce.

\begin{table}[t]
\centering
\caption{Operational state of the 23-model study population under each condition.
RC: reproducibly correct (10/10). NR: non-reproducible. RI: reproducibly incorrect
(0/10). Each row sums to 23.}
\label{tab:conditions}
\begin{tabular}{lrrr}
\toprule
Condition & RC & NR & RI \\
\midrule
\textbf{Task specification} & \textbf{17} & 4 & \textbf{2} \\
Expert role & 2 & 9 & 12 \\
Objective emphasis & 1 & 8 & 14 \\
Anti-hallucination & 1 & 6 & 16 \\
Chain of thought & 1 & 6 & 16 \\
Encouragement & 1 & 5 & 17 \\
Error avoidance & 0 & 7 & 16 \\
Threat & 0 & 7 & 16 \\
\midrule
Untreated & 0 & 7 & 16 \\
\bottomrule
\end{tabular}
\end{table}

\paragraph{Transitions between reproducibly incorrect, non-reproducible, and
reproducibly correct states.}
The quantity of interest is the transition each model makes from its untreated
state, separately for the two failure modes (Table~\ref{tab:transitions}). The task
specification moves thirteen of the sixteen reproducibly incorrect models to
reproducible correctness and four of the seven non-reproducible ones, rescuing
seventeen of 23 in a single pass with no change to the question and no sampling,
voting or retry. Every conventional intervention rescues at most one model from
either group.

An aggregate correctness figure cannot express this separation, and in two
conditions it points the wrong way. Encouragement yields fewer correct executions
than no intervention at all, 30 of 230 against 34, while the number of reproducibly
incorrect models rises from sixteen to seventeen: the aggregate falls, and so does
the visibility of the failure. Chain-of-thought yields more correct executions than
no intervention, 38 of 232 against 34, while leaving the reproducibly incorrect
count unchanged and driving one non-reproducible model into it. In both cases the
aggregate moves in one direction and the operational state of the fleet in the
other.

What the conventional interventions mostly do is move models within and between the
two failure modes --- between inconsistent and consistent error --- which changes an
aggregate figure without changing whether any individual model can be relied on for
the decision. The specification is the only condition that empties most of the
reproducibly incorrect set, and the only one that adds to it nowhere.

\begin{table}[t]
\centering
\caption{Transitions from each untreated state. Left: where the sixteen reproducibly
incorrect models land. Right: where the seven non-reproducible models land. Rescued
is the total reaching reproducible correctness.}
\label{tab:transitions}
\begin{tabular}{lrrrcrrrcr}
\toprule
& \multicolumn{3}{c}{from RI ($n=16$)} && \multicolumn{3}{c}{from NR ($n=7$)} && \\
\cmidrule{2-4}\cmidrule{6-8}
Condition & RC & NR & RI && RC & NR & RI && Rescued \\
\midrule
\textbf{Task specification} & \textbf{13} & 1 & 2 && \textbf{4} & 3 & \textbf{0} && \textbf{17} \\
Expert role & 1 & 3 & 12 && 1 & 6 & 0 && 2 \\
Objective emphasis & 0 & 2 & 14 && 1 & 6 & 0 && 1 \\
Anti-hallucination & 1 & 0 & 15 && 0 & 6 & 1 && 1 \\
Chain of thought & 1 & 0 & 15 && 0 & 6 & 1 && 1 \\
Encouragement & 1 & 0 & 15 && 0 & 5 & 2 && 1 \\
Threat & 0 & 1 & 15 && 0 & 6 & 1 && 0 \\
Error avoidance & 0 & 1 & 15 && 0 & 6 & 1 && 0 \\
\bottomrule
\end{tabular}
\end{table}
